\section{多目标函数与价值分账}
\label{sec:objective}

\subsection{评价边界}

多目标评价在实际功率、状态和产品交付量确定后执行，不反向修改4.2--4.7的物理结果。
运行优化与全生命周期评价分层：4.10用运行期经济目标、可靠性底线和环境代理产生调度；
资产CAPEX、固定运维、融资、替换和税费在给定建设容量的生命周期账簿中另行评价。
因此，运行净收益不等同于项目利润、NPV、IRR或可融资结论。

\subsection{运行经济目标}

定义运行期经济净成本
\begin{equation}
\begin{aligned}
F_{eco}={}&C_{OM}+C_{start}+C_{deg}+C_{transport}
+C_{ENS}+C_{SLA}+C_{curt}+C_{terminal}\\
&-R_{elec}-R_{H2}-R_{compute}-R_{marine}.
\end{aligned}
\label{eq:economic-objective}
\end{equation}
各项统一以CNY计。电力收入按海缆受端交付电量计算：
\begin{equation}
R_{elec}=\sum_t\pi_t^{elec}P_t^{cab,r}\Delta t_t;
\label{eq:revenue-electricity}
\end{equation}
氢气收入按扣除通道损失后的交付量计算：
\begin{equation}
R_{H2}=\sum_t\pi_t^{H2}m_t^{H2,del};
\label{eq:revenue-hydrogen}
\end{equation}
算力和海洋服务收入分别依据4.6的实际服务量和式~\eqref{eq:marine-service}的实际交付量
计算。内部电力转移价只用于业务单元边际贡献，不得在源侧成本和算力成本中重复计入。

$C_{deg}$按电池实际吞吐、SOC和温度相关寿命模型计量；$C_{start}$按实际增机或启动事件
计量；$C_{transport}$按管输/船运实际提取或交付量计量；$C_{terminal}$只用于反映库存
机会成本或滚动终端价值，不能替代物理安全库存约束。所有价格、违约、退化、运输和
终端价值参数须记录币值年、含税口径、适用区域和来源类型，未签认时均为\assumptiontag。

绿氨尚未进入现行物理变量、资产账簿和交付边界，故式~\eqref{eq:economic-objective}不
确认绿氨收入。待4.3补齐合成氨工艺后，应新增可追溯的产量、能耗、库存、运输和合同项。

\subsection{全生命周期经济性账簿}

对给定候选容量，资本回收因子为
\begin{equation}
CRF(r,n)=\frac{r(1+r)^n}{(1+r)^n-1},
\label{eq:crf}
\end{equation}
设备$i$的年化固定成本可写为
\begin{equation}
C_i^{annual}=CAPEX_i\,CRF(r,n_i)+C_i^{FOM}+C_i^{repl,annual}.
\label{eq:annualized-cost}
\end{equation}
该结构与公开技术经济评价中显式计入资本成本、寿命和固定运维的原则一致
\cite{nrel-atb,iec60300}。完整NPV按投运年现金流写为
\begin{equation}
NPV=-I_0+\sum_{y=1}^{Y}
\frac{R_y-C_y-I_y^{repl}}{(1+r)^y}.
\label{eq:npv}
\end{equation}
折现率、建设期支出曲线、税费、融资结构、寿命和更换年份必须由第6章及企业财务边界
签认；未签认时为\assumptiontag。短时场景不得按$8760/H$直接外推为正式NPV、LCOE、
LCOS或LCOH结论。

\subsection{环境目标}

环境目标采用固定边界下的项目排放减基准情景避免排放：
\begin{equation}
F_{env}=GHG_{project}-GHG_{avoided,baseline}.
\label{eq:environment-objective}
\end{equation}
运行代理可按源侧使用、电池吞吐、制氢、运输、氢转电和外购电等活动量乘相应排放因子，
并按电力、氢气、算力和海洋服务的冻结基准计算避免排放。完整生命周期评价必须依据
ISO~14040和ISO~14044完成目标与范围、清单、影响评价和解释
\cite{iso14040,iso14044}。设备BOM/EPD、海工船舶、备件、冷媒、替换、退役回收及多产品
分摊规则在取得企业或第三方LCA数据前均为\assumptiontag；现行代理不得表述为已认证的
产品碳足迹。

\subsection{可靠性目标}

单个确定性时序的未供能为
\begin{equation}
ENS=\sum_t\left(P_t^{ENS,int}+P_t^{ENS,mar}
+P_t^{ENS,comp}\right)\Delta t_t.
\label{eq:ens}
\end{equation}
不同服务等级分别记账，避免用一个ENS掩盖内部关键负荷、海洋服务和基础算力的差异。
只有存在概率场景$s$及其概率$\pi_s$时才定义
\begin{equation}
EENS=\sum_{s\in\mathcal S^{scen}}\pi_sENS_s,
\qquad \sum_s\pi_s=1.
\label{eq:eens}
\end{equation}
EENS定义和场景口径可参见ENTSO-E资源充裕度方法\cite{entsoe-maf}。LOLE、LOLP、N-1、
CVaR、频率最低点、ROCOF和黑启动成功率尚未由小时模型直接计算，不得写作已完成指标。

\subsection{多目标集成}

由于CNY、kgCO$_2$e和MWh不能直接相加，主模型采用$\varepsilon$约束法：
\begin{equation}
\min F_{eco}
\quad\mathrm{s.t.}\quad
ENS\le\varepsilon_{ENS},\quad
F_{env}\le\varepsilon_{GHG},\quad
\eta_{RE}\ge\varepsilon_{RE},
\label{eq:epsilon-constraint}
\end{equation}
其中可再生能源利用率为
\begin{equation}
\eta_{RE}=\frac{\sum_tP_t^{src,use}\Delta t_t}
{\sum_tP_t^{src,av}\Delta t_t}.
\label{eq:renewable-utilization}
\end{equation}
$\varepsilon$约束方法依据Mavrotas的增强实现\cite{mavrotas2009}。阈值表达管理偏好和合同
底线，未经决策者签认时为\assumptiontag。扫描阈值可形成Pareto候选点，但本章只定义
方法，不呈现或宣称最终最优前沿。
